\documentclass[11pt]{article}
\usepackage[margin=1in]{geometry}

\usepackage{amsmath,amssymb}
\usepackage{amsthm}
\usepackage{booktabs}
\usepackage{multirow}
\usepackage{adjustbox}
\usepackage{algorithm}
\usepackage{algpseudocode}
\usepackage{natbib}
\usepackage{pifont}
\newcommand{\cmark}{\ding{51}}
\newcommand{\xmark}{\ding{55}}

\theoremstyle{plain}
\newtheorem{theorem}{Theorem}[section]
\newtheorem{lemma}[theorem]{Lemma}
\newtheorem{corollary}[theorem]{Corollary}
\theoremstyle{remark}
\newtheorem{remark}[theorem]{Remark}

\setlength{\parindent}{0pt}
\setlength{\parskip}{1\baselineskip}

\title{Untitled}
\author{Author Name}
\date{\today}

\begin{document}
\maketitle

% ============================================================
\begin{abstract}
  Residual decoding algorithms score elements, prune one, and update state iteratively.
  Classical algorithms exploit additive statistics such as degree to perform each update
  in $O(n)$, achieving $O(n^2)$ total cost. We show that this efficiency is
  algebraically unavailable to GNNs: exact additive updates require affine score
  encodings, which bounded nonlinear activations cannot implement, and vanishing score
  margins cause approximation errors to invert removal orderings at scale, cascading
  into trajectory divergence. This forces a full network rerun after every deletion,
  inflating cost to $O(n^3)$.

  We show that the objective gradient provides an exact local correction after each
  vertex deletion, maintainable in $O(n)$ per step. Supplying this live gradient signal
  to a lightweight learned updater $U_\theta$ exposes an additive signal whose first
  linear layer can represent the exact residual correction before any nonlinearity fires,
  recovering $O(n^2)$ decoding.

  We validate on \textsc{Max-Clique} using planted-clique instances and a decoder
  ablation spanning the full $O(n^3)$-to-$O(n^2)$ spectrum. Cached-gradient variants
  match or exceed the $O(n^3)$ rerun baseline across all difficulty regimes.
\end{abstract}

% ============================================================
\section{Introduction}
% ============================================================

A wide class of combinatorial algorithms proceeds by a \emph{residual loop}: score
the current elements, prune or select one, update the problem state, and repeat.
This loop underlies greedy peeling~\citep{feige2010}, spectral
rounding~\citep{alon1998}, branch-and-bound solvers, and local greedy methods more
broadly. Its tractability
rests on the \emph{additivity} of the maintained statistics. Degree is the
simplest example: when a vertex is removed, each survivor's degree changes by at
most one, depending only on local adjacency. This local, fixed correction makes the
full update $O(n)$, giving $O(n^2)$ total decoding on dense graphs.

\paragraph{Problem setting.}
We focus on \textsc{Max-Clique}: given a graph, find the largest subset of vertices
that are all pairwise adjacent. The decision version is NP-complete and the
optimization problem is NP-hard, making clique recovery a notorious testbed for
algorithmic heuristics. Classical decoders attack it with greedy top-$k$ selection,
local-degree refinement, spectral rounding, and related residual scoring rules. In
the planted-clique setting used in our experiments, the clique is hidden inside a
random background graph and the decoder must recover it. This makes
\textsc{Max-Clique} a natural place to ask whether learned scorers can replace these
hand-designed residual heuristics without losing their efficient update structure.

\paragraph{GNNs as learned scorers.}
GNNs have emerged as the natural candidate for replacing hand-crafted statistics
with learned scoring rules~\citep{karalias2020, min2022, garmendia2024}. Trained
end-to-end on problem instances, they produce vertex rankings for \textsc{Max-Clique},
\textsc{Max-Cut}, \textsc{Max-Independent-Set}, and related problems without
domain-specific feature engineering. Recent interpretability work shows that GNNs
trained on combinatorial tasks tend to rediscover classical algorithmic
heuristics~\citep{cappart2023}, and for \textsc{Max-Clique} specifically, degree-based
vertex ranking emerges consistently across architectures and training
regimes~\citep{shoham2026ltr}. This raises an optimistic picture: a GNN that has
internalized degree-based scoring should, in principle, slot directly into the
residual loop in place of the classical scorer.

\paragraph{The efficiency gap.}
This hope runs into a precise algebraic obstruction. A score encoding $\varphi$
supports an exact additive update $\varphi(d-1) = \varphi(d) - c$ if and only if
$\varphi$ is affine (Theorem~\ref{thm:affine}). Any nonlinear activation is
incompatible: a bounded function cannot implement an unbounded affine map
(Theorem~\ref{thm:nonlinear}). The problem compounds with scale: as instance size
grows, score gaps between adjacent candidates vanish, so even a fixed approximation
error suffices to invert the removal ordering on arbitrarily large instances
(Theorem~\ref{thm:approx_flip}; Corollary~\ref{cor:vanishing_margin}). Once an
inversion occurs, the approximate and exact decoders evolve on structurally different
residual graphs, and subsequent errors grow without bound rather than averaging out
(Theorem~\ref{thm:propagation}). The only safe remedy is to rerun the GNN after every
deletion---$O(n^2)$ per step, $O(n^3)$ total.

\paragraph{The gradient as an additive coordinate.}
The obstruction above does not rule out efficient decoding; it rules out efficient
decoding from the \emph{raw GNN score} alone. We show that the objective gradient
provides exactly the additive information that is missing. For any locally
decomposable objective, deleting vertex $v$ changes each surviving gradient
coordinate $g_i$ by a quantity determined entirely by the local factor between $i$
and $v$ (Lemma~\ref{lem:local_gradient_correction}). For the clique objective this
correction is an affine function of the deleted vertex's cached probability. This
means the gradient can be updated in $O(n)$ per deletion by maintaining three
sufficient statistics, giving $O(n^2)$ total
maintenance cost after the initial $O(n^2)$ GNN pass
(Theorem~\ref{thm:cached_gradient_complexity}).

Supplying this live cached gradient to a learned score updater $U_\theta$ turns the
problem around: the exact residual correction $\Delta g_i$ is a linear function of
$U_\theta$'s inputs when $g_i$ is available, so the first linear layer of $U_\theta$
can represent it exactly before any nonlinearity fires
(Theorem~\ref{thm:gu_positive}). Without the gradient, recovering $\Delta g_i$ from
the nonlinear score $\sigma(z_{v})$ requires inverting $\sigma$---exactly the
impossibility established in Theorem~\ref{thm:nonlinear}.

\paragraph{Contributions.}
On the theory side we prove:
\begin{itemize}
  \item \textbf{Theorem~\ref{thm:affine}.} Exact additive updates force affine score
    encodings: any score that updates correctly after every deletion must be linear
    in the maintained statistic.
  \item \textbf{Theorem~\ref{thm:nonlinear}.} Nonlinear GNNs cannot be affine:
    bounded activations are incompatible with the unbounded affine structure that
    exact updates require.
  \item \textbf{Theorem~\ref{thm:approx_flip} and Corollary~\ref{cor:vanishing_margin}.}
    Vanishing margins make approximation errors destructive at scale: when score gaps
    shrink with instance size, any fixed-accuracy model will eventually invert the
    removal ordering.
  \item \textbf{Theorem~\ref{thm:propagation}.} Ordering errors compound: a single
    wrong deletion sends the decoder onto a structurally different residual graph,
    after which errors are unbounded.
  \item \textbf{Lemma~\ref{lem:local_gradient_correction} and
    Theorem~\ref{thm:gu_positive}.} The gradient provides an exact local correction:
    deleting a vertex changes each surviving gradient coordinate by a local quantity,
    which a learned updater can represent exactly in its first linear layer.
  \item \textbf{Theorem~\ref{thm:cached_gradient_complexity}.} Cached sufficient
    statistics recover $O(n^2)$ decoding: maintaining $(\mathbf{A}p, \mathbf{B}p,
    \sum p_i)$ with $O(n)$ per-deletion updates gives $O(n^2)$ total cost after the
    initial GNN pass.
\end{itemize}



On the empirical side, we introduce a learned score updater decoder (UPR) that runs
the GNN once and updates scores incrementally via $U_\theta$, and a corresponding
training regime (PR) that supervises $U_\theta$ by imitating a rerun-based teacher.
Both are described in full in Section~\ref{sec:method}. Our findings:
\begin{itemize}
  \item \textbf{Running time.} UPR reduces decoding from $O(n^3)$ to $O(n^2)$
    with no loss in approximation ratio relative to the rerun baseline.
  \item \textbf{Recovery quality.} SGNN+GU under PR training achieves perfect
    approximation and exact recovery on Easy and Medium regimes, matching the
    $O(n^3)$ rerun baseline at $O(n^2)$ cost.
  \item \textbf{Necessity of the gradient channel.} Removing the gradient from
    $U_\theta$ collapses incremental performance even under PR training, confirming
    that the gradient is a necessary architectural condition, not merely a helpful
    signal.
  \item \textbf{Ablation over training regimes and decoders.} Objective-only
    training with the gradient channel already substantially improves over standard
    GNNs; PR training recovers the rerun baseline entirely. Without the gradient,
    PR training makes performance strictly worse than objective training alone.
\end{itemize}

% On the empirical side, we evaluate on planted clique across an ablation of training regimes and decoder processes. The gradient channel is both necessary and sufficient: removing it collapses
% performance even under explicit imitation training, while providing it with
% cached-deletion maintenance recovers the $O(n^3)$ rerun baseline at $O(n^2)$ cost, with the same approximation ratio.

\paragraph{Paper organization.}
Section~\ref{sec:related} reviews related work. Section~\ref{sec:method} develops
the methodology. Section~\ref{sec:theory} presents the theoretical analysis.
Section~\ref{sec:experiments} presents the empirical evaluation.
Section~\ref{sec:discussion} concludes.

% ============================================================
\section{Related Work}
\label{sec:related}
% ============================================================

\paragraph{Message-passing GNNs.}
A message-passing neural network (MPNN) maintains a hidden state
$h_v^{(\ell)} \in \mathbb{R}^d$ updated over $L$ layers as
\[
  h_v^{(\ell+1)}
  =
  \sigma\!\left(
    \mathbf{M}\!\left[
      h_v^{(\ell)},\,
      \mathrm{Agg}_{u\in\mathcal{N}(v)} h_u^{(\ell)}
    \right]
  \right),
\]
with a per-vertex score $z_v = \mathbf{W}_{\mathrm{out}} h_v^{(L)}$ read off at the
end~\citep{xu2019}. Their expressive power is bounded by the 1-WL test~\citep{xu2019,
morris2019}, though our results concern an orthogonal limitation---\emph{update
compatibility}---that persists even beyond the WL bound. Since these networks
ultimately output vertex scores, a natural question is whether those scores resemble
the heuristics used by classical algorithms.

\paragraph{What GNNs learn on combinatorial problems.}
This score-based view connects directly to recent interpretability results.
\citet{cappart2023} argue that GNNs on NP-hard problems behave more like
polynomial-time heuristics than exact solvers, and \citet{shoham2026ltr} give a
concrete \textsc{Max-Clique} example: across architectures and datasets, learned
representations consistently recover degree-based vertex rankings. At first glance,
this is exactly what one would hope for, since degree is the canonical statistic behind
fast residual clique heuristics. The difficulty is that a useful residual heuristic is
not only a good ranking rule; it is also a rule whose state can be updated cheaply
after each deletion. Our work shows that this second property does not automatically
transfer from degree to a nonlinear GNN score.

\paragraph{GNNs for combinatorial optimization.}
The issue becomes central in neural combinatorial optimization, where learned scores
are routinely placed inside residual decoders. \citet{karalias2020} introduced an
unsupervised framework deriving GNN loss directly from the problem objective;
\citet{min2022} extended this with the scattering-based architecture we build on.
Similar unsupervised objectives have been applied to \textsc{Max-Cut},
\textsc{Max-Independent-Set}, and \textsc{TSP}~\citep{schuetz2022,garmendia2024}.
These works show that GNN scores can guide strong combinatorial heuristics, but the
standard decoding pattern still reruns or reuses the network in a residual loop. That
loop is precisely where update compatibility and the vanishing-margin obstruction
(Theorem~\ref{thm:approx_flip}) become critical. This points to a broader design
principle: neural solvers should align not only with the objective they optimize, but
also with the algorithmic update structure used at inference time.

\paragraph{Architectural alignment with algorithms.}
This principle is visible in recent work on algorithm-aligned architectures.
\citet{yau2023} (OptGNN) proved that polynomial-sized MPNNs can implement SDP
gradient descent for Max-CSPs, obtaining UGC-optimal guarantees in theory and
outperforming SDP solvers in practice. Our work is complementary: OptGNN aligns
message passing with the dynamics of a convex relaxation, whereas we align the
residual decoder with the local objective change caused by a deletion. Thus the
algorithmic object we need is not a full relaxation trajectory, but a live marginal
signal that can be updated as the residual graph changes.

\paragraph{Gradient and marginal signals in neural solvers.}
Several recent neural solvers also enrich message passing with dynamic marginal or
algorithmic signals. QRF-GNN~\citep{gaile2024} feeds intermediate GNN predictions
back as dynamic node features to improve convergence within a single inference
pass---structurally related to our gradient channel, but targeting a different goal.
\citet{hadou2024} unroll a dual ascent algorithm into coupled networks so each layer
computes a meaningful algorithmic step. Our approach differs from both: we expose the
objective gradient as a live residual coordinate and let the model learn how to use
it, with the additional structural contribution that the gradient is maintained
exactly in $O(n)$ per deletion via cached sufficient statistics. This leaves a second
question: how should the updater learn to turn this maintained signal into the same
decisions a full rerun would make?

\paragraph{Imitation and self-supervised decoding.}
This training question connects our PR regime to imitation learning in combinatorial
optimization~\citep{liu2021}, where a GNN is trained to replicate a stronger
algorithm's decisions. The novelty here is self-imitation: teacher and student share
weights, and the student must predict the rerun-based removal sequence from
incrementally updated scores using the cached gradient.

% ============================================================
\section{Methodology}
\label{sec:method}
% ============================================================

We evaluate three model families:
\[
  \text{SGNN},\qquad
  \text{SGNN+U},\qquad
  \text{SGNN+GU}.
\]
SGNN is a standard message-passing GNN. SGNN+U adds a learned score updater
$U_\theta$ without gradient access. SGNN+GU provides $U_\theta$ with the current
cached objective gradient. The central question is whether the gradient channel is
necessary, sufficient, or both for efficient incremental decoding.

\subsection{Problem Setting}
\label{sec:problem}

We evaluate on planted clique on $G(n,\tfrac{1}{2},k)$ random graphs. To sample
the Erd\H{o}s--R\'enyi background, each unordered vertex pair is included
independently with probability $1/2$. We then choose $k$ vertices uniformly at random
and add all missing edges among them, forming the planted clique $C^\star$. The goal
is to recover a large valid clique, ideally $C^\star$. Following~\citet{shoham2026ltr},
we partition instances by known theoretical thresholds:
\begin{itemize}
  \item \textbf{Easy} ($k\in[62,100]$, $n=1000$): degree-based heuristics
    succeed~\citep{kucera1995}.
  \item \textbf{Medium} ($k\in[36,61]$): requires iterative residual
    methods~\citep{alon1998, feige2010}.
  \item \textbf{Hard} ($k\in[20,35]$): no polynomial-time algorithm is known to
    succeed~\citep{manurangsi2021}.
\end{itemize}

\paragraph{Metrics.}
We evaluate solution quality and representation structure through five complementary
measures.

\textit{Approximation ratio} $= |C_{\text{found}}|/|C^\star|$ measures the size of
the recovered clique relative to the planted clique. A ratio of $1.0$ means the
decoder found a clique at least as large as the planted one, though not necessarily
identical to it.

\textit{Overlap} $= |C_{\text{found}} \cap C^\star|/|C^\star|$ measures what
fraction of the planted clique vertices were actually recovered. Unlike approximation
ratio, overlap is sensitive to whether the correct vertices were found, not just
whether a large clique was found.

\textit{Exact recovery} $= \mathbf{1}[C_{\text{found}} = C^\star]$ is the strictest
criterion: the recovered clique must be identical to the planted one.

\textit{Explained variance} of the first two principal components (PC1, PC2) of the
final GNN hidden layer $h_v^{(L)}$, computed via PCA across all vertices. Following
\citet{shoham2026ltr}, PC1 typically captures over $95\%$ of the variance,
revealing whether the learned representation collapses onto a dominant scalar
ordering direction.

\textit{Spearman correlation} $\rho$ between degree and the PC1 projection of
$h_v^{(L)}$, and between degree and the output scores $s_v$. A high $|\rho|$
indicates that the model has learned a degree-based ranking, consistent with the
classical LPR heuristic~\citep{shoham2026ltr}. We also report $\rho$ between the
post-update scores $\tilde{s}_v$ (after one application of $U_\theta$) and degree,
to measure whether the gradient update preserves or disrupts the learned ranking.

\subsection{Model Architecture}
\label{sec:architectures}

All models share the same message-passing backbone: $L=4$ shared message-passing
steps, hidden dimension $64$, and $\sigma=\tanh$. The input feature is a scalar
liveness indicator; no degree, clique-membership, or structural feature is provided.

\paragraph{Standard GNN (SGNN).}
The input feature is first projected as $h_v^{(0)} = W_{\mathrm{in}} x_v + b_{\mathrm{in}}$,
where $x_v \in \mathbb{R}$ is a scalar liveness indicator initialized to $1$. At layer $\ell$, SGNN
aggregates neighbor representations and updates:
\[
  q_v^{(\ell)} = \frac{1}{n-1}\sum_{u\in\mathcal{N}(v)} h_u^{(\ell)},
  \qquad
  h_v^{(\ell+1)} = \sigma\!\left(M_{\mathrm{upd}}\!\left[h_v^{(\ell)},
  q_v^{(\ell)}\right]\right),
\]
\[
  s_v = W_{\mathrm{out}} h_v^{(L)} + b_{\mathrm{out}},
  \qquad
  p_v = \operatorname{sigmoid}(s_v),
\]
where $M_{\mathrm{upd}}$ is a two-layer MLP with hidden dimension $d$, $\sigma$
is a nonlinear activation applied after each layer, $W_{\mathrm{out}}, b_{\mathrm{out}}$
project the final hidden state to a scalar score after $L$ message-passing steps,
and $p_v \in (0,1)$ is the vertex probability used by the clique objective.


\paragraph{SGNN+U and SGNN+GU.}
Both use the SGNN backbone and add a lightweight MLP updater $U_\theta$. After
removing vertex $v^*$, each surviving score is updated as
\[
  s_i \leftarrow s_i + U_\theta(s_i,\, s_{v^*},\, A_{iv^*},\, \gamma_i,\, \gamma_{v^*}),
\]
where $s_i$ and $s_{v^*}$ are the current scores of vertex $i$ and the removed
vertex, $A_{iv^*}$ is the adjacency indicator between them, and $\gamma_i,
\gamma_{v^*}$ are gradient channels that differ between the two variants.
For SGNN+U, $\gamma_i = \gamma_{v^*} = 0$: the updater receives no gradient
information. For SGNN+GU,
\[
  \gamma_i = g_i = \frac{\partial \mathcal{L}_{\mathrm{clique}}}{\partial p_i},
  \qquad
  \gamma_{v^*} = g_{v^*} = \frac{\partial \mathcal{L}_{\mathrm{clique}}}{\partial p_{v^*}},
\]
where $g_i$ and $g_{v^*}$ are maintained via the cached sufficient statistics
described in Section~\ref{sec:cached_maintenance}.

\subsection{Cached Gradient Maintenance}
\label{sec:cached_maintenance}

Recomputing the full gradient after every deletion costs $O(n^2)$ per step and
$O(n^3)$ total. For the clique objective, this is unnecessary. The gradient
\[
  g_i = -\frac{2(\mathbf{A}p)_i}{k(k-1)}
  + \frac{2(\mathbf{B}p)_i}{k(n-k)}
  + \frac{20(\sum_j p_j - k)}{n^2}
\]
depends only on three sufficient statistics $q^A = \mathbf{A}p$, $q^B = \mathbf{B}p$,
$S = \sum_j p_j$. After deleting vertex $v$ with cached probability $p_v$, each
statistic is updated by a single rank-one subtraction:
\[
  q^A \leftarrow q^A - \mathbf{A}_{:v}\,p_v,
  \qquad
  q^B \leftarrow q^B - \mathbf{B}_{:v}\,p_v,
  \qquad
  S \leftarrow S - p_v.
\]
This costs $O(n)$ per deletion. The cache is initialized in $O(n^2)$ alongside
the initial GNN pass, so the total decoding cost remains $O(n^2)$
(Theorem~\ref{thm:cached_gradient_complexity}).

\subsection{Decoding Procedures}
\label{sec:decoding}

All decoders run the GNN once, then iteratively remove the lowest-scoring vertex
until the remaining set forms a clique, followed by an add-back pass that restores
any removed vertex compatible with the found clique. They differ only in how---and
whether---scores are updated after each deletion.

\paragraph{Top-$k$ ($O(n^2)$).}
The GNN runs once. Vertices are ranked by score in descending order and greedily
added to the solution: a vertex is included if it is adjacent to all previously
selected vertices, stopping when $k$ vertices are found~\citep{shoham2026ltr}.
No pruning loop is used.

\paragraph{Rerun-Pruned ($O(n^3)$).}
The full GNN is rerun on the current residual graph after every deletion to produce
fresh scores. The lowest-scoring vertex is removed at each step until the remaining
set forms a clique, followed by an add-back pass. This is the only $O(n^3)$ decoder
and serves as the adaptive upper bound.

\paragraph{LPR ($O(n^2)$).}
The GNN runs once. Scores are never updated; the lowest-scoring vertex is removed
at each step until a clique remains, followed by an add-back pass~\citep{shoham2026ltr}.

\paragraph{UPR ($O(n^2)$).}
The GNN runs once. After each deletion of $v^*$, $U_\theta$ updates all surviving
scores. For SGNN+U, $U_\theta$ receives no gradient. For SGNN+GU, $U_\theta$
receives the cached deletion-updated gradient, and the score update is gated to
neighbors of $v^*$ only.

\begin{algorithm}[t]
  \caption{UPR: Cached-Deletion Decoder}
  \label{alg:gub}
  \begin{algorithmic}[1]
    \State Run GNN once to obtain scores $s_i$ and probabilities $p_i$ for all $i$.
    \State Build gradient cache from $\mathcal{L}$ and $p$. \hfill\Comment{$O(n^2)$}
    \State $C \leftarrow V(G)$;\; $H \leftarrow []$.
    \While{$C$ does not form a clique}
    \State $v^* \leftarrow \arg\min_{i \in C} s_i$;\; append $v^*$ to $H$;\; $C \leftarrow C \setminus \{v^*\}$.
    \State Remove $v^*$'s contribution from gradient cache. \hfill\Comment{$O(n)$, by Lemma~\ref{lem:local_gradient_correction}}
    \State $g_i \leftarrow$ gradient from updated cache for all $i \in C$. \hfill\Comment{$O(n)$}
    \State $s_i \mathrel{+}= U_\theta(s_i,\, s_{v^*},\, A_{iv^*},\, g_i,\, g_{v^*})$ for all $i \in C$. \hfill\Comment{$O(n)$}
    \EndWhile
    \For{$v$ from last to first in $H$}
    \If{$C \cup \{v\}$ is a clique} $C \leftarrow C \cup \{v\}$. \EndIf
    \EndFor
    \State \Return $C$.
  \end{algorithmic}
\end{algorithm}

\subsection{Training}
\label{sec:training}

All models are trained only on Medium-difficulty instances. Each run uses 64
epochs, with 1000 newly sampled instances per epoch, minibatch size 32, and no
overlap with the held-out evaluation instances. The unsupervised clique objective is
\[
  \mathcal{L}_{\mathrm{clique}}(p; A, k)
  =
  -\frac{p^\top \mathbf{A} p}{k(k-1)}
  +
  \frac{p^\top \mathbf{B} p}{k(n-k)}
  +
  10\!\left(\frac{\textstyle\sum_i p_i - k}{n}\right)^{\!2},
\]
where $p_i = \sigma(s_i)$ and $\mathbf{B}$ is the complement adjacency matrix.

\paragraph{Objective training (OBJ).}
$\mathcal{L}_{\mathrm{train}} = \mathcal{L}_{\mathrm{clique}}$. $U_\theta$ is
supervised only through the objective, with no direct signal on the removal sequence.

\paragraph{Probable-Remover training (PR).}
PR adds an imitation loss that directly supervises the incremental decoder. At each
step $t$, a rerun-based teacher reruns the GNN on the current residual graph $G_t$
and selects the lowest-scoring vertex $v^t_{\mathrm{teach}}$. A student---running
the GNN only once and updating scores incrementally via $U_\theta$---is penalised
for disagreeing:
\[
  \mathcal{L}_{\mathrm{PR}} = \frac{1}{T} \sum_{t=0}^{T-1}
  \mathrm{CE}_{\mathrm{softmax}}\!\left(-\tilde{s}^{(t)}_{V(G_t)},\;
  v^t_{\mathrm{teach}}\right),
  \qquad
  \mathcal{L}_{\mathrm{train}} = \mathcal{L}_{\mathrm{clique}} +
  \mathcal{L}_{\mathrm{PR}}.
\]
For SGNN+U, $U_\theta$ receives no gradient signal; it is trained purely through
backpropagation of $\mathcal{L}_{\mathrm{PR}}$ with $\gamma_i = \gamma_{v^*} = 0$.
For SGNN+GU, the cached gradient is updated in $O(n)$ after each teacher removal
before being passed to $U_\theta$. The teacher cost is $O(n^3)$ per training
instance, paid at train time only.

\begin{algorithm}[t]
  \caption{PR Training Step}
  \label{alg:pr_training}
  \begin{algorithmic}[1]
    \State Sample instance $(G, k)$; run GNN once to obtain $\tilde{s}$, $p$.
    \State Build gradient cache from $\mathcal{L}$ and $p$. \hfill\Comment{$O(n^2)$; SGNN+GU only}
    \For{$t = 0, \ldots, T-1$}
    \State $v^t_{\mathrm{teach}} \leftarrow \arg\min_{i \in V(G_t)} \mathrm{GNN}(G_t)_i$. \hfill\Comment{Teacher reruns GNN, $O(n^2)$}
    \State $\mathcal{L}_{\mathrm{PR}} \mathrel{+}= \mathrm{CE}_{\mathrm{softmax}}(-\tilde{s}_{V(G_t)},\, v^t_{\mathrm{teach}})$. \hfill\Comment{Student predicts same removal}
    \State $G_t \leftarrow G_t \setminus \{v^t_{\mathrm{teach}}\}$.
    \State Update gradient cache by removing $v^t_{\mathrm{teach}}$. \hfill\Comment{$O(n)$; SGNN+GU only}
    \State $\tilde{s}_i \mathrel{+}= U_\theta(\tilde{s}_i,\, \tilde{s}_{v^t},\, A_{iv^t},\, \gamma_i,\, \gamma_{v^t})$ for all $i \in V(G_t)$.
    \EndFor
    \State Optimise $\mathcal{L}_{\mathrm{clique}} + \mathcal{L}_{\mathrm{PR}} / T$.
  \end{algorithmic}
\end{algorithm}

% ============================================================
% ============================================================
\section{Theoretical Analysis}
\label{sec:theory}
% ============================================================

We develop the theory in three acts. The first shows why nonlinear GNN scores are
structurally incompatible with efficient incremental decoding. The second identifies
the objective gradient as the missing additive coordinate. The third shows how to
maintain that coordinate in $O(n)$ per deletion.

\paragraph{Act I: Why nonlinear scores fail.}
The efficiency of classical residual algorithms rests on a simple algebraic fact: if
removing a neighbor decreases degree by exactly $1$, and the score is $s_i =
\varphi(d_i)$, then an exact additive update requires $\varphi(d-1) = \varphi(d) - c$
for a fixed constant $c$. This forces $\varphi$ to be affine.

\begin{theorem}[Affine encoding is necessary for exact additive updates]
  \label{thm:affine}
  Let $d_i \in \mathbb{Z}_{\geq 0}$ be an additive statistic decreasing by $1$ when
  a neighbor is deleted. If the score $s_i = \varphi(d_i)$ satisfies
  $\varphi(d-1) = \varphi(d) - c$ for a fixed constant $c$ and all relevant $d$,
  then $\varphi(d) = cd + b$ for some $b \in \mathbb{R}$.
\end{theorem}
\begin{proof}[Proof sketch]
  The condition pins every first discrete difference of $\varphi$ to the same
  constant $c$, which forces $\varphi$ to be linear. Full proof in
  Appendix~\ref{app:proofs}.
\end{proof}


GNNs with nonlinear activations cannot implement this: the score after a deletion
passes through $\sigma$ on a different input than the score before, and since
$\sigma$ is nonlinear, $\sigma(a - b) \neq \sigma(a) - \sigma(b)$ in general.
No fixed additive correction can bridge this gap.

\begin{theorem}[Nonlinear GNNs cannot implement exact additive updates]
  \label{thm:nonlinear}
  Let $f_\theta\colon\mathcal{G}\to\mathbb{R}^n$ be an MPNN with learned parameters
  $\theta$, whose update MLPs contain at least one nonlinear activation $\sigma$. Then for any fixed $c$,
  there exist $G$ and adjacent $u,v$ such that $f_\theta(G)_v - c \neq
  f_\theta(G\setminus\{u\})_v$ regardless of $\theta$.
\end{theorem}
\begin{proof}
  By Theorem~\ref{thm:affine}, exact updates require $f_\theta(G)_v$ to be affine
  in $d_v$. But $f_\theta(G)_v$ is a composition of $\sigma$ with linear
  functions of the neighborhood aggregation. Since $\sigma$ is nonlinear,
  $\sigma(a - b) \neq \sigma(a) - \sigma(b)$ in general: the score after a
  deletion is not the score before minus a fixed constant. No choice of $\theta$
  can eliminate this gap, because it is a consequence of the nonlinearity of
  $\sigma$ itself, not of the parameter values.
\end{proof}


% \begin{remark}
%   The impossibility is also statistical: any strictly increasing reparametrisation
%   $\varphi(d_i)$ satisfies the same pairwise ranking constraints as $d_i$, so
%   ranking supervision cannot enforce the additive scale that exact updates require.
% \end{remark}

The failure is not merely local. As instance size grows, score gaps between adjacent
candidates vanish. A fixed approximation error then suffices to invert the removal
ordering, and a single inversion sends the two decoders onto structurally different
residual graphs from which there is no recovery.









\begin{theorem}[Approximate scores induce ordering errors at scale]
  \label{thm:approx_flip}
  Let $\hat{s}_i^{(n)} = s_i^{(n)} + \varepsilon_i^{(n)}$ and let $\gamma_n :=
  \min_{i\neq j}|s_i^{(n)} - s_j^{(n)}|$ be the minimum nonzero score gap. If
  each removal step introduces a fixed nonzero error $\varepsilon > 0$, then after
  $r = \lceil \gamma_n / \varepsilon \rceil$ steps the accumulated error crosses
  $\gamma_n$ and a pairwise ordering is inverted.
\end{theorem}
\begin{proof}[Proof sketch]
  Each removal adds $\varepsilon$ to the cumulative score error. After $r$ steps
  the total error is $r\varepsilon$. Once $r\varepsilon > \gamma_n$, the error
  exceeds the closest score gap and the wrong vertex is removed. Since $\varepsilon
  > 0$ is fixed and the trajectory length grows with $n$, this crossing is
  inevitable. Full proof in Appendix~\ref{app:proofs}.
\end{proof}

\begin{corollary}[Any fixed approximation error eventually causes a flip]
  \label{cor:vanishing_margin}
  For any nonlinear GNN score with irreducible per-step approximation error,
  the accumulated error will eventually exceed the minimum score gap $\gamma_n$,
  inverting the removal ordering. Exact additive coordinates are the only
  structural remedy: they update the score algebraically and incur no approximation
  error per step.
\end{corollary}
\begin{proof}[Proof sketch]
  Immediate from Theorem~\ref{thm:approx_flip}: with fixed $\varepsilon > 0$ and
  trajectory length growing with $n$, the accumulated error $r\varepsilon$
  eventually exceeds $\gamma_n$ for any instance large enough.
\end{proof}


\begin{theorem}[Local ordering errors cascade into trajectory divergence]
  \label{thm:propagation}
  If at step $t$ estimated scores invert the ordering of $i$ and $j$, the two
  decoders remove different vertices. All subsequent score discrepancies equal
  differences between GNN outputs on structurally different graphs and can be
  arbitrarily large.
\end{theorem}
\begin{proof}[Proof sketch]
  After the inversion, vertex $i$ is present in one residual graph and absent from
  the other. Every future score is computed on a different input; the deviation is
  no longer a perturbation of the original error vector but an uncontrolled
  structural difference. Full proof in Appendix~\ref{app:proofs}.
\end{proof}




Together, Theorems~\ref{thm:affine}--\ref{thm:propagation} establish that no
nonlinear GNN can support cheap incremental decoding from its raw score alone, and
that the damage compounds rather than dissipates.

\paragraph{Act II: The gradient as an additive coordinate.}
The obstruction rules out efficient decoding from the raw score, not from a
richer signal. For any locally decomposable objective, deleting a vertex removes
exactly its local contribution from each surviving gradient coordinate---a
correction determined entirely by the local factor between the deleted vertex and
each survivor.

\begin{lemma}[Local deletion gives a local gradient correction]
  \label{lem:local_gradient_correction}
  Let \[\mathcal{L}(p;G) = \sum_i \phi_i(p_i) + \sum_{(i,j)\in F(G)}
  \psi_{ij}(p_i,p_j,A_{ij})\] and $g_i(G) = \partial\mathcal{L}/\partial p_i$.
  After deleting vertex $v$,
  \[
    g_i(G\setminus\{v\}) - g_i(G)
    =
    -\frac{\partial \psi_{iv}(p_i,p_v,A_{iv})}{\partial p_i},
  \]
  with the right-hand side zero when $(i,v)\notin F(G)$.
\end{lemma}
\begin{proof}[Proof sketch]
  Differentiating the sum, every term except $\psi_{iv}$ survives deletion of $v$
  unchanged. The correction is therefore exactly the removed local derivative.
  Full proof in Appendix~\ref{app:proofs}.
\end{proof}

For quadratic objectives this correction is an affine function of $(A_{iv}, p_v)$,
both of which are available to a learned updater. This is the key representability
result.

\begin{theorem}[Gradient-updater representability]
  \label{thm:gu_positive}
  For quadratic objectives $\mathcal{L}=\sum_r\alpha_r p^\top M_r p$, the
  deletion correction is
  \[
    \Delta g_i(v) = -2\sum_r\alpha_r(M_r)_{iv}\,p_v,
  \]
  a linear function of $(A_{iv}, p_v)$. An updater $U_\theta$ receiving gradient
  information can represent $\Delta g_i$ exactly in its first linear layer, before
  any nonlinear activation fires.
\end{theorem}
\begin{proof}[Proof sketch]
  The correction is linear in the inputs already available to $U_\theta$; no
  nonlinear inversion of $\sigma$ is required. Without $g_i$, recovering $\Delta
  g_i$ from $\sigma(z_v)$ requires inverting $\sigma$, which
  Theorem~\ref{thm:nonlinear} shows is impossible with a fixed additive rule. Full
  proof in Appendix~\ref{app:proofs}.
\end{proof}

\paragraph{Act III: Cached maintenance recovers $O(n^2)$.}
The gradient correction is cheap to maintain. For the clique objective, $g_i$
depends only on three sufficient statistics $(\mathbf{A}p, \mathbf{B}p, \sum p_j)$.
Each deletion subtracts the deleted vertex's column contribution from these
statistics in $O(n)$, and the updated gradient is read off in $O(n)$ with no
matrix operations.

\begin{theorem}[Cached gradient maintenance yields $O(n^2)$ decoding]
  \label{thm:cached_gradient_complexity}
  The sufficient statistics $(\mathbf{A}p, \mathbf{B}p, \sum_i p_i)$ are
  initialized in $O(n^2)$. Each deletion updates them in $O(n)$ by subtracting the
  deleted vertex's column contribution. Over $n$ deletions the total update cost is
  $O(n^2)$, matching the initial GNN pass.
\end{theorem}
\begin{proof}[Proof sketch]
  Initialization requires two matrix-vector products at $O(n^2)$ each. Each
  deletion subtracts two length-$n$ vectors and one scalar at $O(n)$. Summing over
  at most $n$ deletions gives $O(n^2)$. Full proof in Appendix~\ref{app:proofs}.
\end{proof}

\begin{remark}
  The cache tracks gradient changes due to vertex deletions only, not score drift
  from $U_\theta$. It should be understood as a deletion-updated residual signal,
  not as the exact gradient of the current score vector after each learned update.
\end{remark}

The explicit clique-gradient correction used at each step is:
\begin{equation}
  \label{eq:clique_correction}
  \Delta g_i(v)
  =
  \frac{2A_{iv}p_v}{k(k-1)}
  -
  \frac{2B_{iv}p_v}{k(n-k)}
  -
  \frac{20p_v}{n^2},
\end{equation}
read in $O(1)$ per entry from the updated cache after the $O(n)$ column subtraction.




% ============================================================
\section{Experiments}
\label{sec:experiments}
% ============================================================

\subsection{Setup}
We use the architectures from Section~\ref{sec:architectures} and the training
protocol from Section~\ref{sec:training}. All experiments use planted-clique
instances with $n=1000$; evaluation set sizes are reported in the table captions.
Unless otherwise stated, reported values are mean $\pm$ standard deviation over
the corresponding held-out instances. Hardware: AMD Ryzen 9 5950X, 64\,GB RAM, NVIDIA RTX 3080;
Python 3.11, PyTorch 2.2.

\subsection{Main Results}
\label{sec:main_results}

Table~\ref{tab:main} reports planted clique recovery across all model and decoder
combinations.

\begin{table}[t]
  \centering
  \caption{Planted clique recovery on $n=1000$ across Easy, Medium, and Hard
    regimes, reported as mean$\,\pm\,$std over 10 held-out instances per regime;
    entries without $\pm$ are exactly $0.00$ or $1.00$ across all instances.
    \textbf{Bold} marks the highlighted best results: SGNN+GU+PR with
    UPR is the strongest $O(n^2)$ incremental decoder on Easy and Medium, while
    SGNN+U+PR with Rerun is best on Hard. Runtime reports asymptotic decoding cost:
    all decoders are $O(n^2)$ except Rerun, which is $O(n^3)$. SGNN has no updater and
  cannot run UPR.}
  \label{tab:main}
  \resizebox{\textwidth}{!}{%
    \begin{tabular}{lllccccccccc}
      \toprule
      & & & \multicolumn{3}{c}{Easy} & \multicolumn{3}{c}{Medium}
      & \multicolumn{3}{c}{Hard} \\
      \cmidrule(lr){4-6}\cmidrule(lr){7-9}\cmidrule(lr){10-12}
      Model & Decoder & Runtime & Approx & Overlap & Exact
      & Approx & Overlap & Exact
      & Approx & Overlap & Exact \\
      \midrule
      \multirow{3}{*}{SGNN + OBJ}
      & Top-$k$   & $O(n^2)$ & $1.00$ & $1.00$ & $1.00$
      & $0.75{\pm}0.43$ & $0.74{\pm}0.45$ & $0.67{\pm}0.58$
      & $0.46{\pm}0.12$ & $0.08{\pm}0.07$ & $0.00$ \\
      & Rerun     & $O(n^3)$ & $1.00$ & $1.00$ & $1.00$
      & $1.00$ & $1.00$ & $1.00$
      & $0.49{\pm}0.09$ & $0.00$ & $0.00$ \\
      & LPR       & $O(n^2)$ & $1.00$ & $1.00$ & $1.00$
      & $0.75{\pm}0.43$ & $0.74{\pm}0.45$ & $0.67{\pm}0.58$
      & $0.46{\pm}0.12$ & $0.08{\pm}0.07$ & $0.00$ \\
      \midrule
      \multirow{4}{*}{SGNN+U + OBJ}
      & Top-$k$   & $O(n^2)$ & $1.00$ & $1.00$ & $1.00$
      & $0.62{\pm}0.36$ & $0.58{\pm}0.40$ & $0.33{\pm}0.58$
      & $0.37{\pm}0.06$ & $0.08{\pm}0.07$ & $0.00$ \\
      & Rerun     & $O(n^3)$ & $1.00$ & $1.00$ & $1.00$
      & $1.00$ & $1.00$ & $1.00$
      & $0.63{\pm}0.33$ & $0.36{\pm}0.56$ & $0.33{\pm}0.58$ \\
      & LPR       & $O(n^2)$ & $1.00$ & $1.00$ & $1.00$
      & $0.62{\pm}0.36$ & $0.58{\pm}0.40$ & $0.33{\pm}0.58$
      & $0.37{\pm}0.06$ & $0.08{\pm}0.07$ & $0.00$ \\
      & UPR       & $O(n^2)$ & $0.13{\pm}0.01$ & $0.00$ & $0.00$
      & $0.17{\pm}0.01$ & $0.00$ & $0.00$
      & $0.34{\pm}0.09$ & $0.00$ & $0.00$ \\
      \midrule
      \multirow{4}{*}{SGNN+GU + OBJ}
      & Top-$k$   & $O(n^2)$ & $1.00$ & $1.00$ & $1.00$
      & $1.00$ & $1.00$ & $1.00$
      & $0.48{\pm}0.06$ & $0.06{\pm}0.03$ & $0.00$ \\
      & Rerun     & $O(n^3)$ & $1.00$ & $1.00$ & $1.00$
      & $1.00$ & $1.00$ & $1.00$
      & $0.54{\pm}0.08$ & $0.06{\pm}0.02$ & $0.00$ \\
      & LPR       & $O(n^2)$ & $1.00$ & $1.00$ & $1.00$
      & $1.00$ & $1.00$ & $1.00$
      & $0.50{\pm}0.05$ & $0.05{\pm}0.05$ & $0.00$ \\
      & UPR       & $O(n^2)$ & $0.12{\pm}0.00$ & $0.00{\pm}0.01$ & $0.00$
      & $0.18{\pm}0.02$ & $0.00$ & $0.00$
      & $0.42{\pm}0.04$ & $0.00$ & $0.00$ \\
      \midrule
      \multirow{4}{*}{SGNN+U + PR}
      & Top-$k$   & $O(n^2)$ & $1.00$ & $1.00$ & $1.00$
      & $0.55{\pm}0.40$ & $0.48{\pm}0.47$ & $0.33{\pm}0.58$
      & $0.34{\pm}0.11$ & $0.04{\pm}0.01$ & $0.00$ \\
      & Rerun     & $O(n^3)$ & $1.00$ & $1.00$ & $1.00$
      & $1.00$ & $1.00$ & $1.00$
      & $\mathbf{0.82{\pm}0.31}$ & $\mathbf{0.67{\pm}0.58}$ & $\mathbf{0.67{\pm}0.58}$ \\
      & LPR       & $O(n^2)$ & $1.00$ & $1.00$ & $1.00$
      & $0.55{\pm}0.40$ & $0.48{\pm}0.47$ & $0.33{\pm}0.58$
      & $0.34{\pm}0.11$ & $0.04{\pm}0.01$ & $0.00$ \\
      & UPR       & $O(n^2)$ & $0.11{\pm}0.01$ & $0.00$ & $0.00$
      & $0.18{\pm}0.03$ & $0.00$ & $0.00$
      & $0.29{\pm}0.04$ & $0.00$ & $0.00$ \\
      \midrule
      \multirow{4}{*}{SGNN+GU + PR}
      & Top-$k$   & $O(n^2)$ & $1.00$ & $1.00$ & $1.00$
      & $1.00$ & $1.00$ & $1.00$
      & $0.47{\pm}0.05$ & $0.02{\pm}0.03$ & $0.00$ \\
      & Rerun     & $O(n^3)$ & $1.00$ & $1.00$ & $1.00$
      & $1.00$ & $1.00$ & $1.00$
      & $0.50{\pm}0.07$ & $0.03{\pm}0.03$ & $0.00$ \\
      & LPR       & $O(n^2)$ & $1.00$ & $1.00$ & $1.00$
      & $0.46{\pm}0.21$ & $0.37{\pm}0.33$ & $0.00$
      & $0.47{\pm}0.05$ & $0.02{\pm}0.03$ & $0.00$ \\
      & UPR       & $O(n^2)$ & $\mathbf{1.00}$ & $\mathbf{1.00}$ & $\mathbf{1.00}$
      & $\mathbf{1.00}$ & $\mathbf{1.00}$ & $\mathbf{1.00}$
      & $0.48{\pm}0.09$ & $0.04{\pm}0.04$ & $0.00$ \\
      \bottomrule
    \end{tabular}%
  }
\end{table}


\paragraph{Rerun decoding is strong but expensive.}
Rerun decoding gives perfect Easy and Medium recovery for every model,
confirming that a strong GNN ranking suffices when fresh scores are available after
every deletion~\citep{shoham2026ltr}. Its cost is consistently higher:
Table~\ref{tab:timing} reports $0.76$--$0.78$\,s per instance for rerun,
compared with $0.20$--$0.21$\,s for LPR and $0.37$--$0.57$\,s for UPR decoder.


\paragraph{Without the gradient channel, UPR collapses.}
SGNN+U can still rank vertices reasonably under Top-$k$ and LPR, but its updater
does not work as a residual decoder. UPR stays near zero exact recovery: on Easy
it reaches only $0.11-0.13{\pm}0.01$ under differet training regimes, while on
Medium it remains at $0.00$. Only Rerun restores perfect Easy/Medium recovery,
at the cost of recomputing the GNN after every deletion. PR supervision alone is
therefore not enough without the gradient channel, and can even hurt UPR decoder.


\paragraph{The gradient channel is necessary and PR makes it usable.}
SGNN+GU under objective-only training still performs poorly under UPR decoder
($0.12-0.18{\pm}0.02$ on Easy/Medium), showing that the
gradient channel alone is not enough---the model must also be trained to use it
via PR training. Adding PR training transforms UPR decoding to perfect approximation,
overlap, and exact recovery on Easy and Medium. The gradient channel makes the PR
training objective achievable; PR training drives the model to exploit the gradient
signal at decoding time.

%  Probably remove!
\paragraph{PR training trades static ranking for update quality.}
For SGNN+GU, PR training hurts LPR ($1.00\to0.46$ on Medium) while enabling
perfect UPR recovery ($1.00$). The model sacrifices one-pass score quality to
learn an effective incremental update. For SGNN+U this trade-off does not occur:
without the gradient channel, PR training cannot teach $U_\theta$ a useful update
rule, so both LPR and UPR remain poor.


\paragraph{Hard regime.}
Hard instances remain difficult for all models, consistent with the planted clique
hardness conjecture~\citep{manurangsi2021}. The strongest Hard result overall is
SGNN+U+PR with rerun decoding, which reaches $0.82{\pm}0.31$ approximation but
requires $O(n^3)$ time. Among $O(n^2)$ decoders, Hard is not solved reliably:
SGNN+GU+OBJ with LPR has the highest approximation at $0.50{\pm}0.05$, while
SGNN+GU+PR with UPR is the strongest UPR variant at $0.48{\pm}0.09$.

\begin{table}[t]
  \centering
  \caption{Wall-clock time in seconds per instance on planted-clique graphs with
    $n=1000$, reported as mean $\pm$ standard deviation over 1000 held-out instances
  per regime and averaged over the Easy, Medium, and Hard regimes.}
  \label{tab:timing}
  \begin{tabular}{lcccc}
    \toprule
    Model & Top-$k$ & Rerun & LPR & UPR \\
    \midrule
    SGNN + OBJ      & 0.06 & 0.78 & 0.21 & --- \\
    SGNN+U + OBJ    & 0.06 & 0.77 & 0.20 & 0.38 \\
    SGNN+GU + OBJ   & 0.06 & 0.78 & 0.21 & 0.57 \\
    SGNN+U + PR     & 0.06 & 0.76 & 0.20 & 0.37 \\
    SGNN+GU + PR    & 0.06 & 0.76 & 0.20 & 0.54 \\
    \bottomrule
  \end{tabular}
\end{table}


\subsection{Embedding Geometry}
\label{sec:pca}

Table~\ref{tab:pca} reports PCA diagnostics and Spearman correlations on the final
GNN hidden layer and scalar outputs.

\begin{table}[t]
  \centering
  \caption{Representation diagnostics on 1000 held-out test instances per regime
    ($n=1000$), reported as $|\rho|\pm\sigma$ (mean absolute Spearman $\pm$ std,
    all $p=0$ to machine precision unless noted).
    PC1/PC2: variance explained by the first two principal components of $h_v^{(L)}$.
    $|\rho|(\text{PC1},\text{deg})$ and $|\rho|(\text{logits},\text{deg})$: absolute
    Spearman between PC1 / output logits and residual degree.
    These four quantities are identical across all models, training regimes, and
    difficulty regimes: PC1$=1.00$, PC2$=0.00$,
    $|\rho|(\text{PC1},\text{deg})=1.00\pm0.00$,
    $|\rho|(\text{logits},\text{deg})=1.00\pm0.00$.
    $|\rho|(\tilde{s},\text{deg})$: absolute Spearman between post-update scores and
    residual degree after one application of $U_\theta$ (E$=$Easy, M$=$Medium,
    H$=$Hard). Despite similar global $|\rho|$ values across models at step~1,
    only SGNN+GU+PR avoids removing above-median-degree vertices during UPR:
    \emph{Above-med.}\ reports the fraction of UPR removals where the removed
    vertex has residual degree strictly above the alive median (medium regime,
    $n{=}10$ graphs $\times$ 500 steps, pooled over all steps; k-less autograd).
  SGNN has no updater and cannot run UPR.}
  \label{tab:pca}
  \resizebox{\textwidth}{!}{%
    \begin{tabular}{llcccccccccccccc}
      \toprule
      & & \multicolumn{2}{c}{PC var.} &
      \multicolumn{3}{c}{$|\rho|(\text{PC1},\text{deg})$} &
      \multicolumn{3}{c}{$|\rho|(\text{logits},\text{deg})$} &
      \multicolumn{3}{c}{$|\rho|(\tilde{s},\text{deg})$} &
      \multicolumn{2}{c}{Above-med.\ (\%)} \\
      \cmidrule(lr){3-4}\cmidrule(lr){5-7}\cmidrule(lr){8-10}\cmidrule(lr){11-13}\cmidrule(lr){14-15}
      Model & Train & PC1 & PC2 & E & M & H & E & M & H & E & M & H & Pooled & Mean$\pm$std \\
      \midrule
      SGNN & OBJ
      & 1.00 & 0.00
      & $1.00\pm0.00$ & $1.00\pm0.00$ & $1.00\pm0.00$
      & $1.00\pm0.00$ & $1.00\pm0.00$ & $1.00\pm0.00$
      & \multicolumn{3}{c}{--- (no updater)}
      & \multicolumn{2}{c}{---} \\
      \midrule
      \multirow{2}{*}{SGNN+U}
      & OBJ
      & 1.00 & 0.00
      & $1.00\pm0.00$ & $1.00\pm0.00$ & $1.00\pm0.00$
      & $1.00\pm0.00$ & $1.00\pm0.00$ & $1.00\pm0.00$
      & $0.52\pm0.01$ & $0.52\pm0.01$ & $0.50\pm0.01$
      & --- & --- \\
      & PR
      & 1.00 & 0.00
      & $1.00\pm0.00$ & $1.00\pm0.00$ & $1.00\pm0.00$
      & $1.00\pm0.00$ & $1.00\pm0.00$ & $1.00\pm0.00$
      & $0.50\pm0.04$ & $0.50\pm0.03$ & $0.49\pm0.02$
      & $49.7$ & $49.7\pm2.5$ \\
      \midrule
      \multirow{2}{*}{SGNN+GU}
      & OBJ
      & 1.00 & 0.00
      & $1.00\pm0.00$ & $1.00\pm0.00$ & $1.00\pm0.00$
      & $1.00\pm0.00$ & $1.00\pm0.00$ & $1.00\pm0.00$
      & $0.51\pm0.01$ & $0.47\pm0.03$ & $0.36\pm0.17$
      & --- & --- \\
      & PR
      & 1.00 & 0.00
      & $1.00\pm0.00$ & $1.00\pm0.00$ & $1.00\pm0.00$
      & $1.00\pm0.00$ & $1.00\pm0.00$ & $1.00\pm0.00$
      & $0.49\pm0.03$ & $0.49\pm0.04$ & $0.51\pm0.02$
      & $\mathbf{0.0}$ & $\mathbf{0.0\pm0.0}$ \\
      \bottomrule
  \end{tabular}}
\end{table}


\paragraph{The GNN output is perfectly degree-ranked.}
The strongest signal is already present before any updater is applied: the final
GNN output logits have absolute Spearman correlation $1.00{\pm}0.00$ with vertex
degree across all models and regimes. Thus the backbone output learns a perfect
degree ranking, up to orientation, regardless of architecture or training regime.
The embeddings show the same collapse: PC1 explains essentially all variance
(PC1$=1.00$, PC2$=0.00$), matching the degree-collapse finding
of~\citet{shoham2026ltr}.

% TODO add table for this
\paragraph{The update disrupts degree alignment.}
After one application of $U_\theta$, the absolute Spearman correlation between
residual degree and updated scores drops to $0.36$--$0.52$ across updater variants
and regimes. This is not necessarily a failure: degree ranking is useful for the
initial GNN score, but the update is meant to track residual removal information,
not preserve raw degree order. SGNN+GU+PR achieves perfect Easy/Medium UPR recovery
despite only moderate degree alignment after the update, indicating that the
successful updater is using gradient-informed information not captured by degree
alone.


% \paragraph{SGNN+U and SGNN+GU embeddings are identical in structure.}
% Both have strictly one-dimensional embeddings (PC2$=0.00$), unlike what a richer
% gradient-informed backbone might produce. The gradient channel in SGNN+GU enters
% only through $U_\theta$ at decoding time, not through the GNN backbone, so the
% backbone representations are structurally the same. The difference between SGNN+U
% and SGNN+GU lies in whether $U_\theta$ receives the live gradient, and this
% difference determines whether PR can train UPR to succeed rather than collapse.

\section{Discussion}
\label{sec:discussion}

The experiments point to a distinction that is easy to obscure when evaluating
learned combinatorial decoders: learning a good initial ranking is not the same
as learning a ranking that remains valid under residual updates. In our setting,
the GNN backbone reliably learns a degree-like ordering, and rerunning the GNN
after each deletion recovers strong Easy and Medium performance. Yet this does
not by itself yield an efficient decoder. The difficulty is not that the initial
score is uninformative; it is that the score is not intrinsically compatible with
the residual operation applied by the decoder.

This explains why the updater ablations are decisive. Without the gradient
channel, PR training gives the updater the right target sequence but not the
right state variable. The model is asked to imitate a rerun-based decoder from
a nonlinear score whose residual change is not algebraically available. This
is why SGNN+U fails under UPR even when trained with PR. In contrast, SGNN+GU
exposes a live gradient-derived coordinate whose deletion effect is local and
maintainable. PR training then has a usable object to learn from: it does not
need to invent the residual correction from a static nonlinear score, but only
to learn how to use the maintained gradient signal to match the rerun decisions.

The results also clarify the role of cached maintenance. The cache is not merely
an implementation trick for speeding up the same computation. It changes what
information is available to the decoder throughout the pruning trajectory. A
stale gradient loses alignment because it ignores the vertices already removed;
a deletion-updated cache preserves the dominant residual change at $O(n)$ cost
per step. Thus the $O(n^2)$ decoder is not obtained by approximating an
$O(n^3)$ rerun more cheaply at each step, but by maintaining the right residual
quantity incrementally.

% ============================================================
% KNOWN-k ASSUMPTION:
% Experiments use the planted clique size k as the stopping target.
% If k is unknown, wrap decoding in binary search over k; O(log n) runs.
% ============================================================
\paragraph{Known target size.}
Our experiments assume that the target clique size $k$ is known, as in the
planted-clique setup. The decoder therefore knows when to stop pruning. If $k$ is
not known, the same procedure can be wrapped in a binary search over candidate
target sizes, with the returned set checked for clique validity. This requires
$O(\log n)$ decoder calls, so an $O(n^2)$ UPR decoder becomes
$O(n^2\log n)$ rather than $O(n^2)$.

A useful consequence is that the successful updater need not preserve raw degree
alignment. The representation diagnostics show that the GNN score itself is
degree-aligned, while the post-update scores are only moderately correlated with
degree. This is consistent with the purpose of UPR: after deletion begins, the
decoder should track residual removal information, not simply reproduce the
initial degree ordering. The gradient channel supplies information beyond this
static ranking, which explains why SGNN+GU+PR can recover Easy and Medium
instances even when the updated score is no longer a pure degree proxy. This also
leaves an open interpretability question: if the successful updater is not simply
following degree, what concept has it learned to use for each residual deletion
decision? Our diagnostics identify the gradient channel as necessary, but they do
not fully characterize the internal rule learned by $U_\theta$.

More broadly, the paper suggests that neural decoders should be judged not only
by score quality but by update compatibility. Classical residual algorithms are
efficient because their maintained statistics have explicit local update rules.
Our results show that a learned decoder can recover this property by exposing
and maintaining a gradient coordinate. The problem-specific part is the form of
the objective gradient and its cached sufficient statistics; the algorithmic
principle is to align the learned decoder with the residual algebra of the
objective.

\paragraph{Limitations and future work.}
The theory shows that gradient injection makes additive residual coordinates
representable, but it does not guarantee that SGD will discover and use such
coordinates consistently. GU addresses this by keeping the gradient-derived
signal explicit during decoding. A stronger solution would be an architecture
or loss that makes the GNN itself learn update-compatible coordinates, ideally
removing the need for a separate updater. Another limitation is that GU assumes
the objective and its derivative are known in advance, so the required cached
statistics can be derived before decoding. Extending this idea to settings where
the relevant loss or gradient is unknown, learned, or only available as a
black-box signal remains open. Our training procedure also relies on PR
supervision from a rerun-based teacher, which is expensive even though inference
is quadratic. Reducing this training cost is an important practical direction.
Finally, the experiments focus on planted clique. Extending the same
gradient-maintenance principle to other locally decomposable objectives, such
as independent set, cut, and constraint-satisfaction relaxations, remains an
important test of the generality of the approach. In particular, it remains open
how to obtain an analogous $O(n^2)$ learned residual decoder for
\textsc{Max-Cut}, where the residual decision structure differs from clique
peeling.


\bibliographystyle{plainnat}
\bibliography{references}

% ============================================================
\appendix
\section{Gradient Rectification}
\label{app:rectification}
% ============================================================

\begin{theorem}[Gradient in $U_\theta$ enables exact additive updates]
  \label{thm:rectification}
  Let $\mathcal{L}(p) = \sum_r\alpha_r p^\top M_r p$ and
  $g_i = \partial\mathcal{L}/\partial p_i$. The correct residual correction after
  deleting $v^*$ is $\Delta g_i = -2\sum_r\alpha_r(M_r)_{iv^*}p_{v^*}$. Then:
  \begin{enumerate}
    \item \textbf{With $g_i$ in $U_\theta$:} $\Delta g_i$ is linear in
      $(A_{iv^*}, p_{v^*})$, inputs already available to $U_\theta$. The first linear
      layer of $U_\theta$ can represent $\Delta g_i$ exactly before any nonlinearity
      fires.
    \item \textbf{Without $g_i$ in $U_\theta$:} $U_\theta$ receives
      $s_{v^*}=\sigma(z_{v^*})$, a nonlinear encoding. Recovering $\Delta g_i$
      requires inverting $\sigma$, which by Theorem~\ref{thm:nonlinear} no fixed
      additive rule can achieve. Errors accumulate and compound by
      Theorem~\ref{thm:propagation}.
  \end{enumerate}
\end{theorem}
\begin{proof}
  \textbf{Part 1.} The correction $\Delta g_i = -2\sum_r\alpha_r(M_r)_{iv^*}p_{v^*}$
  is linear in $(A_{iv^*}, p_{v^*})$. The first layer of $U_\theta$ computes
  $W[s_i, s_{v^*}, A_{iv^*}, g_i, g_{v^*}]$, a linear map that can select the
  appropriate combination of $g_{v^*}$ and $A_{iv^*}$ to represent $\Delta g_i$
  exactly before any activation is applied.

  \textbf{Part 2.} Without $g_i$, recovering $\Delta g_i$ from $\sigma(z_{v^*})$
  requires inverting $\sigma$. By Theorem~\ref{thm:affine}, the correct update
  requires affine dependence on the underlying statistic $d_{v^*}$; by
  Theorem~\ref{thm:nonlinear}, no fixed additive rule achieves this. Every step
  incurs nonzero error, and Theorem~\ref{thm:propagation} guarantees these errors
  compound into trajectory divergence.
\end{proof}

\begin{remark}
  Injecting the gradient into the GNN backbone rather than into $U_\theta$ does not
  help: the backbone gradient is fixed at the end of the initial forward pass and goes
  stale after the first deletion. $U_\theta$ requires the \emph{current} gradient,
  which is supplied by the cached sufficient statistics updated in $O(n)$ at each step.
  This is why SGNN+GU outperforms MAGNN+U empirically.
\end{remark}











\appendix

% ============================================================
\section{Full Proofs}
\label{app:proofs}
% ============================================================

We give complete proofs of all results stated in Section~\ref{sec:theory}, following
the structure of the proof document that accompanies this paper.

\subsection*{Act I: The algebraic obstruction}

\begin{theorem}[Theorem~\ref{thm:affine} restated]
  Let $\mathcal{X}$ be a state space and $\varphi:\mathcal{X}\to\mathbb{R}^p$ a
  score encoding. Suppose that for every elementary residual transformation
  $T_\alpha$ there exists a fixed vector $c_\alpha\in\mathbb{R}^p$ such that
  $\varphi(T_\alpha x)=\varphi(x)+c_\alpha$ for all admissible states $x$. Then,
  for every state obtained from a reference state $x_0$ by a sequence
  $T_{\alpha_r}\cdots T_{\alpha_1}$,
  \[
    \varphi(T_{\alpha_r}\cdots T_{\alpha_1}x_0)
    =
    \varphi(x_0)+\sum_{q=1}^r c_{\alpha_q}.
  \]
  In particular, if the residual state is represented by coordinates
  $d\in\mathbb{Z}^m$ and elementary transformations induce fixed increments
  $\Delta_\alpha$, then $\varphi$ is affine on the generated lattice:
  $\varphi(d)=Cd+b$.
\end{theorem}
\begin{proof}
  The first identity follows by repeated substitution of the assumed update rule:
  \begin{align*}
    \varphi(T_{\alpha_r}\cdots T_{\alpha_1}x_0)
    &= \varphi(T_{\alpha_{r-1}}\cdots T_{\alpha_1}x_0) + c_{\alpha_r} \\
    &= \cdots = \varphi(x_0) + \sum_{q=1}^r c_{\alpha_q}.
  \end{align*}
  In coordinate form, the relation $\varphi(d+\Delta_\alpha)-\varphi(d)=c_\alpha$
  says that all discrete first differences of $\varphi$ along the generated
  directions are constant. Hence $\varphi$ is affine on each connected lattice
  component generated by these directions, giving $\varphi(d) = Cd + b$ for
  $b = \varphi(0)$.
\end{proof}

\begin{theorem}[Theorem~\ref{thm:nonlinear} restated]
  Let $f_\theta\colon\mathcal{G}\to\mathbb{R}^n$ be an MPNN with learned parameters
  $\theta$, whose update MLPs contain at least one nonlinear activation $\sigma$.
  Then for any fixed $c$, there exist $G$ and adjacent $u,v$ such that
  $f_\theta(G)_v - c \neq f_\theta(G\setminus\{u\})_v$ regardless of $\theta$.
\end{theorem}
\begin{proof}
  By Theorem~\ref{thm:affine}, exact updates require $f_\theta(G)_v$ to be affine
  in $d_v$. But $f_\theta(G)_v$ is a composition of $\sigma$ with linear functions
  of the neighborhood aggregation. Since $\sigma$ is nonlinear,
  $\sigma(a - b) \neq \sigma(a) - \sigma(b)$ in general: the score after a deletion
  is not the score before minus a fixed constant. No choice of $\theta$ can
  eliminate this gap, because it is a consequence of the nonlinearity of $\sigma$
  itself, not of the parameter values.
\end{proof}

\begin{theorem}[Theorem~\ref{thm:approx_flip} restated]
  Let $\hat{s}_i^{(n)} = s_i^{(n)} + \varepsilon_i^{(n)}$ and let $\gamma_n :=
  \min_{i\neq j}|s_i^{(n)} - s_j^{(n)}|$ be the minimum nonzero score gap. If
  each removal step introduces a fixed nonzero error $\varepsilon > 0$, then after
  $r = \lceil \gamma_n / \varepsilon \rceil$ steps the accumulated error crosses
  $\gamma_n$ and a pairwise ordering is inverted.
\end{theorem}
\begin{proof}
  Each removal adds $\varepsilon$ to the cumulative score error. After $r$ steps
  the total error is $r\varepsilon$. Once $r\varepsilon > \gamma_n$, the error
  exceeds the closest score gap and the wrong vertex is removed. Since $\varepsilon
  > 0$ is fixed and the trajectory length grows with $n$, this crossing is
  inevitable for any instance large enough.
\end{proof}

\begin{corollary}[Corollary~\ref{cor:vanishing_margin} restated]
  For any nonlinear GNN score with irreducible per-step approximation error,
  the accumulated error will eventually exceed the minimum score gap $\gamma_n$,
  inverting the removal ordering. Exact additive coordinates are the only
  structural remedy: they update the score algebraically and incur no approximation
  error per step.
\end{corollary}
\begin{proof}
  Immediate from Theorem~\ref{thm:approx_flip}: with fixed $\varepsilon > 0$ and
  trajectory length growing with $n$, the accumulated error $r\varepsilon$
  eventually exceeds $\gamma_n$ for any instance large enough.
\end{proof}

\begin{theorem}[Theorem~\ref{thm:propagation} restated]
  Consider an iterative decoder that selects the element with minimum score. If at
  some step approximate scores invert the ordering of two candidate elements, then
  the approximate and exact decoders delete different elements and subsequently
  evolve on different residual states. Later score discrepancies are governed by
  structural divergence of the states, not by the original local approximation
  error.
\end{theorem}
\begin{proof}
  Once the ordering is inverted, the two decoders choose different elements. The
  residual states after the deletion are consequently distinct: one contains vertex
  $i$ and not $j$; the other contains $j$ and not $i$. Every later score is
  evaluated on a different residual structure, so the next discrepancy is a
  difference between two scoring problems on different graphs, not a perturbation of
  the same score vector. Such discrepancies can compound arbitrarily across the
  trajectory.
\end{proof}

\subsection*{Act II: The gradient as an additive coordinate}

\begin{lemma}[Lemma~\ref{lem:local_gradient_correction} restated]
  Let $\mathcal{L}(p;G) = \sum_i \phi_i(p_i) + \sum_{(i,j)\in
  F(G)}\psi_{ij}(p_i,p_j,A_{ij})$ be a locally decomposable objective and
  $g_i(G) = \partial\mathcal{L}/\partial p_i$. After deleting vertex $v$,
  \[
    g_i(G\setminus\{v\}) - g_i(G)
    =
    -\frac{\partial \psi_{iv}(p_i,p_v,A_{iv})}{\partial p_i},
  \]
  with the right-hand side zero when $(i,v)\notin F(G)$.
\end{lemma}
\begin{proof}
  Differentiating the decomposed objective,
  \[
    g_i(G) = \phi_i'(p_i)
    + \sum_{j:(i,j)\in F(G)} \frac{\partial \psi_{ij}(p_i,p_j,A_{ij})}{\partial p_i}.
  \]
  Deleting $v$ removes exactly the factor $\psi_{iv}$ from this sum and leaves all
  other local factors involving $i$ unchanged. The stated correction follows
  immediately.
\end{proof}

\begin{theorem}[Theorem~\ref{thm:gu_positive} restated]
  For quadratic objectives $\mathcal{L}=\sum_r\alpha_r p^\top M_r p$, the deletion
  correction is
  \[
    \Delta g_i(v) = -2\sum_r\alpha_r(M_r)_{iv}\,p_v.
  \]
  This is a linear function of $(A_{iv}, p_v)$. An updater $U_\theta$ receiving
  gradient information can represent $\Delta g_i$ exactly in its first linear layer,
  before any nonlinear activation fires.
\end{theorem}
\begin{proof}
  For a quadratic term $p^\top M_r p$, the $i$-th gradient coordinate is
  $\nabla_i(p^\top M_r p) = 2(M_r p)_i$. Removing $v$ eliminates the $v$-th
  contribution from this marginal, giving
  \[
    \Delta g_i(v) = -2\sum_r \alpha_r (M_r)_{iv} p_v.
  \]
  This is linear in $(A_{iv}, p_v)$, both available as inputs to $U_\theta$. The
  first affine layer of $U_\theta$ computes $W[s_i, s_{v^*}, A_{iv^*}, g_i,
  g_{v^*}]$, a linear map that can set appropriate weights on $(A_{iv}, p_v)$ to
  represent $\Delta g_i$ exactly before any activation is applied.

  Without $g_i$, $U_\theta$ receives $s_{v^*} = \sigma(z_{v^*})$, a nonlinear
  encoding. Recovering $\Delta g_i$ from $\sigma(z_{v^*})$ requires inverting
  $\sigma$. By Theorem~\ref{thm:affine}, the correct update requires affine
  dependence on the underlying statistic; by Theorem~\ref{thm:nonlinear}, no fixed
  additive rule achieves this inversion. Every step therefore incurs nonzero error,
  and Theorem~\ref{thm:propagation} guarantees these errors compound into
  trajectory divergence.
\end{proof}

\subsection*{Act III: Cached maintenance}

\begin{theorem}[Theorem~\ref{thm:cached_gradient_complexity} restated]
  On a dense graph, suppose one GNN forward pass costs $O(n^2)$. The initial
  gradient cache $(\mathbf{A}p, \mathbf{B}p, \sum_i p_i)$ can be built in $O(n^2)$
  by two matrix-vector products. If each deletion updates this cache by subtracting
  the deleted vertex's column contribution, then each residual step costs $O(n)$
  and the full decoding trajectory costs $O(n^2)$.
\end{theorem}
\begin{proof}
  The initial forward pass and the initial products $\mathbf{A}p$ and $\mathbf{B}p$
  each cost $O(n^2)$. At a deletion step, updating $\mathbf{A}p$ and $\mathbf{B}p$
  requires subtracting two length-$n$ vectors $\mathbf{A}_{:v}p_v$ and
  $\mathbf{B}_{:v}p_v$, and updating $S$ requires one scalar subtraction, costing
  $O(n)$ total. Since there are at most $n$ deletions, the total update cost is
  $O(n^2)$, matching the initial setup. Hence the full decoding trajectory costs
  $O(n^2)$.
\end{proof}

\begin{remark}
  The cache maintains the gradient induced by deleting vertices from the current
  residual objective under the cached probability values. If the learned updater
  also changes all probabilities after each deletion, recomputing the exact fresh
  gradient of those changed probabilities would require a dense matrix-vector
  product. The cached gradient should therefore be understood as a
  deletion-updated residual signal used by $U_\theta$, not as a full recomputation
  after arbitrary global score drift.
\end{remark}



\end{document}
